\chapter{Canonical Form Reduction}\label{ch:canonical}

This chapter reduces a general MPS tensor to block-diagonal data via iterated
invariant-subspace decomposition. Two reduction schemes are relevant.
The block-injective reduction aims at injective TP-normalized blocks and the
canonical form conditions of Chapter~\ref{ch:block_perm}. The normal-canonical
reduction uses a weighted family of irreducible TP-normalized blocks with
primitive transfer maps and pairwise distinct nonzero weight moduli, in the
sense of \cite{Cirac2017MPDO_AnnPhys}.
Thus the reductions recorded here are conditional
primitive-block results. They should not be read as the full
translation-invariant canonical-form theorem of
\cite[Theorem~\texttt{Th:TIcanonical}, lines~742--763]{PerezGarcia2007Matrix}.
The normalization convention is also dual to the one used in the cited theorem:
the blocks are written in the unital orientation
$\sum_i A_k^i(A_k^i)^\dagger=\Id$ for
$\E_{A_k}(X)=\sum_i A_k^iX(A_k^i)^\dagger$
\cite[Theorem~\texttt{Th:TIcanonical}, lines~754--758]{PerezGarcia2007Matrix},
whereas this chapter uses
the trace-preserving, or left-canonical, orientation
$\sum_i (A_k^i)^\dagger A_k^i=\Id$.
Passing to the conjugate-transposed Kraus family $K_i := (A^i)^\dagger$
exchanges the two orientations: the transfer map of $K$ is the Frobenius adjoint
of $\E_A$ (formalized in
\texttt{MPS/CanonicalForm/BlockingViaAdjoint.lean}).
The Perron--Frobenius eigenvector existence and TP-gauge construction
are in Chapter~\ref{ch:qpf}; this chapter uses them to produce
left-canonical normalization data and the normal canonical form data.
The reduction draws on
\cite{PerezGarcia2007Matrix}, \cite{Cirac2021Matrix},
and~\cite[Chapter~6]{Wolf2012Quantum}.

\section{Normal tensors, canonical forms, and bases of normal tensors}\label{sec:paper_faithful_defs}

This section records the unstrengthened definitions of \emph{normal tensor},
\emph{canonical form}, and \emph{basis of normal tensors} from
\cite{Cirac2017MPDO_AnnPhys}.  The stronger predicates used elsewhere in
this chapter (Definition~\ref{def:is_normal_canonical_form} and the
canonical-form BNT predicate of Chapter~\ref{ch:bnt}) bundle additional
structure---left-canonical normalization, strict modulus ordering,
one-copy-per-sector---which is not present in the literal source-paper
definitions stated below.

\begin{definition}[CPSV normal tensor (NT)]\label{def:nt_paper}
	\lean{MPSTensor.IsNormalTensor}
	\leanok
	A tensor $A : \texttt{MPSTensor } d\, D$ is a \emph{normal tensor}
	if (i)~$A$ admits no nontrivial invariant orthogonal projection
	(\texttt{IsIrreducibleTensor}), and (ii)~its associated CPM
	$\mathcal{E}_A(X) = \sum_i A^i X (A^i)^\dagger$ has a unique
	eigenvalue of magnitude equal to its spectral radius which is
	equal to one (\texttt{IsPrimitive (transferMap A)}, after
	the spectral-radius normalization in
	\cite[paragraph preceding the NT definition]{Cirac2017MPDO_AnnPhys}).
\end{definition}

\begin{definition}[CPSV canonical form (CF)]\label{def:cf_paper}
	\lean{MPSTensor.IsCanonicalFormPaper}
	\leanok
	A tensor $A$ is in \emph{canonical form} if there is a finite family
	of normal tensors $A_k : \texttt{MPSTensor } d\, D_k$ and weights
	$\mu_k \in \mathbb{C}$ ($k = 1, \dots, r$) with
	\[
		A^i \;\cong\; \bigoplus_{k=1}^r \mu_k\, A_k^i
	\]
	in the sense that $A$ and the block-assembled tensor
	$\texttt{toTensorFromBlocks}\,\mu\,A_\bullet$ generate the same
	MPV family at every length, and the weights satisfy the
	normalization $|\mu_k| \le 1$ for all $k$ and $|\mu_k| = 1$ for
	at least one $k$.
\end{definition}

\begin{definition}[CPSV basis of normal tensors (BNT)]\label{def:bnt_paper}
	\lean{MPSTensor.IsBNTPaper}
	\leanok
	A finite family $\{A_j\}_{j=1}^g$ of normal tensors is a
	\emph{basis of normal tensors} for $A$ if:
	\begin{enumerate}
		\item each $A_j$ is normal in the sense of Definition~\ref{def:nt_paper};
		\item for every $N \ge 1$ there exist coefficients
		      $c_j^{(N)} \in \mathbb{C}$ with
		      \[
		        V^{(N)}(A) \;=\; \sum_{j=1}^{g} c_j^{(N)}\, V^{(N)}(A_j);
		      \]
		\item there exists $N_0$ such that for all $N > N_0$ the family
		      $\bigl(V^{(N)}(A_j)\bigr)_{j=1}^{g}$ is linearly independent.
	\end{enumerate}
\end{definition}

\begin{remark}[Relation to the strong predicates]
	A tensor $A$ which is both irreducible (Definition~\ref{def:irreducible_tensor})
	and has primitive transfer map is normal in the sense of
	Definition~\ref{def:nt_paper}.  The reverse implications from the
	stronger predicates of Definition~\ref{def:is_normal_canonical_form} and
	of Chapter~\ref{ch:bnt} require either the fundamental-theorem step or
	the algebraic-to-spectral normality equivalence; these appear in
	downstream chapters.
\end{remark}

\section{Transfer map normalization}\label{sec:transfer_map_normalization}

\begin{remark}[Scalar normalization conventions]\leanok
	The reductions in this chapter use four routine normalization facts:
	scaling preserves injectivity, the transfer map under a scalar~$\zeta$
	rescales by $|\zeta|^2$, unit-modulus phase rescaling preserves the
	left-canonical condition $\sum_i (A^i)^\dagger A^i = \Id$, and the MPV of a
	weighted block-diagonal tensor factorizes as
	$\mu_k^N V^{(N)}(A_k)_\sigma = \eta_k^N |\mu_k|^N V^{(N)}(A_k)_\sigma$ with
	$\eta_k := \mu_k/|\mu_k|$.  These are scalar identities, formalized in the
	Lean module \texttt{SharedInfra/Scaling.lean}.

	Without first placing each block in the left-canonical normalization, the
	implication
	$\mc{V}(\bigoplus_k \mu_k A_k) = \mc{V}(\bigoplus_k \mu_k B_k)
	\Rightarrow \forall k,\ \mc{V}(A_k)=\mc{V}(B_k)$
	fails even for two scalar blocks with distinct nonzero weights:
	for $\mu=(1,2)$ and $A=(1,3/2)$, $B=(3,1/2)$, the two direct sums generate
	identical MPV families but the first blocks differ at $N=1$.  The
	canonical-form hypotheses
	(Definition~\ref{def:cf_paper})
	remove this ambiguity by combining left-canonical normalization with the
	strict ordering $|\mu_1|>\cdots>|\mu_r|$.
\end{remark}

\section{Invariant subspace decomposition}

\begin{theorem}[Unitary conjugation preserves the MPV family]\label{thm:samempv_unitary}
	\lean{MPSTensor.sameMPV_conj_unitary}
	\leanok
	\uses{def:same_mpv}
	If $U$ is a unitary matrix, then $A$ and $U^\dagger A U$
	generate the same MPV family.
\end{theorem}

\begin{proof}\leanok
	By cyclicity of the trace:
	$\tr((U^\dagger A^{i_1} U) \cdots (U^\dagger A^{i_N} U))
		= \tr(A^{i_1} \cdots A^{i_N})$.
\end{proof}

\begin{definition}[Support projection]\label{def:support_proj}
	\lean{MPSTensor.supportProj}
	\leanok
	For a positive semi-definite matrix $\rho \ge 0$,
	the \emph{support projection} $P$ is the orthogonal
	projection onto the range of $\rho$.
	It is constructed via the spectral decomposition:
	if $\rho = U \diag(\lambda_1,\ldots,\lambda_D) U^\dagger$,
	then $P = U \diag(\mathbf{1}_{\lambda_1>0},\ldots,\mathbf{1}_{\lambda_D>0}) U^\dagger$,
	where $\mathbf{1}_{\lambda_j>0}$ is $1$ if $\lambda_j > 0$
	and $0$ otherwise.
\end{definition}

\begin{definition}[Invariant projection predicate]\label{def:has_inv_proj}
	\lean{MPSTensor.HasInvariantProj}
	\leanok
	\uses{def:mps_tensor}
	An MPS tensor $A$ \emph{has an invariant projection}
	if there exists an orthogonal projection $P$ with
	$P \neq 0$, $P \neq \Id$, and
	$(\Id - P) A^i P = 0$ for all~$i$.
	A tensor is \emph{irreducible} if it has no nontrivial invariant projection
	(Definition~\ref{def:irreducible_tensor}).
\end{definition}

\begin{theorem}[PSD fixed point gives invariant projection]\label{thm:fp_inv_proj}
	\lean{MPSTensor.lowerZero_of_posSemidef_fixedPoint}
	\leanok
	\uses{def:transfer_map}
	Let $\rho \ge 0$ be a fixed point of $\E_A$,
	and let $P$ be the support projection of $\rho$.
	Then for all~$i$, $(\Id - P) A^i P = 0$:
	the range of $P$ is invariant under all $A^i$.
\end{theorem}

\begin{proof}
	\leanok
	From $\E_A(\rho) = \rho$ and positivity of each
	summand $A^i \rho (A^i)^\dagger$,
	$(\Id - P) A^i \rho (A^i)^\dagger (\Id - P) = 0$
	for each~$i$.
	Since $\rho$ is supported on the range of $P$,
	this forces $(\Id - P) A^i P = 0$.

	This is the finite-dimensional version of
	\cite[Proposition~6.10]{Wolf2012Quantum}
	(``Stationary subspaces'': the support projection
	of any stationary state is sub-harmonic, i.e.\
	invariant under the channel).
	The equivalence between irreducibility and
	absence of nontrivial invariant projections is
	\cite[Theorem~6.2]{Wolf2012Quantum};
	see also
	\cite[Section~IV.A.1]{Cirac2021Matrix}.
\end{proof}

\begin{theorem}[Two-block decomposition]\label{thm:two_block_decomp}
	\lean{MPSTensor.exists_twoBlock_decomp_of_posSemidef_fixedPoint_strict}
	\leanok
	\uses{def:transfer_map, def:same_mpv2}
	If the transfer map has a PSD fixed point $\rho$ that is
	\emph{not} positive definite (i.e.\ $\rho$ has a nontrivial
	kernel), then there exist MPS tensors $A_1, A_2$ of smaller
	bond dimensions such that the two-block tensor
	$A_1 \oplus A_2$ generates the same MPV family as~$A$:
	$\mc{V}(A) = \mc{V}(A_1 \oplus A_2)$.
\end{theorem}

\begin{proof}
	\leanok
	\uses{thm:fp_inv_proj, thm:samempv_unitary}
	The support projection $P$ of $\rho$ is neither $0$ nor
	$\Id$ (since $\rho \neq 0$ and $\rho$ is not PD).
	Theorem~\ref{thm:fp_inv_proj} gives
	$(\Id - P) A^i P = 0$ for all~$i$.
	A unitary change of basis that block-diagonalizes $P$ puts
	$A^i$ in upper-triangular block form, with diagonal blocks
	$B_{11}^i$ and $B_{22}^i$.
	Theorem~\ref{thm:samempv_unitary} ensures the MPV is preserved by the
	conjugation.
	The trace of the upper-triangular product equals the sum
	of the diagonal-block traces, so the block-diagonal tensor
	obtained by discarding the off-diagonal blocks has the same MPV
	coefficients as the conjugated tensor, and hence as~$A$.
\end{proof}

\section{Iterated reduction}

\begin{definition}[Irreducible tensor]\label{def:irreducible_tensor}
	\lean{MPSTensor.IsIrreducibleTensor}
	\leanok
	\uses{def:has_inv_proj}
	An MPS tensor~$A$ is \emph{irreducible}
	if it has no nontrivial invariant projection.
\end{definition}

\begin{theorem}[Irreducible block decomposition]\label{thm:irred_decomp}
	\lean{MPSTensor.exists_irreducible_blockDecomp}
	\leanok
	\uses{def:irreducible_tensor}
	Every MPS tensor $A$ admits a block-diagonal tensor
	$\bigoplus_{k=1}^r A_k$
	where each $A_k$ is irreducible and
	$\mc{V}(A) = \mc{V}\!\bigl(\bigoplus_k A_k\bigr)$.
\end{theorem}

\begin{proof}
	\leanok
	\uses{thm:two_block_decomp, thm:pf_eigenvector_existence}
	By strong induction on the bond dimension~$D$.
	If $A$ is already irreducible, take $r = 1$.
	Otherwise, by Definition~\ref{def:irreducible_tensor},
	$A$ has a nontrivial invariant projection~$P$.
	Diagonalizing~$P$ yields the same two-block upper-triangular
	splitting used in Theorem~\ref{thm:two_block_decomp}; the
	block-diagonal tensor obtained by discarding the off-diagonal part
	has the same MPV family as~$A$, and the two diagonal blocks have
	strictly smaller bond dimensions.
	Iterating this splitting produces a block-diagonal tensor
	$\bigoplus_k A_k$ whose blocks are irreducible.
	Apply the induction hypothesis to each block and concatenate.
	(If one prefers a positive-eigenvector argument without assuming TP
	normalization, the correct replacement for the channel-specific
	Ces\`aro theorem is Theorem~\ref{thm:pf_eigenvector_existence};
	the induction itself needs only the invariant projection.)

	Both \cite[Theorem~4]{PerezGarcia2007Matrix} and
	\cite[Section~IV.A.1]{Cirac2021Matrix} obtain the block decomposition by
	the same invariant-projection splitting used here.
	An alternative derivation, developed
	in~\cite[Theorem~6.14]{Wolf2012Quantum},
	characterises the fixed-point set of a quantum channel
	as a $*$-algebra and then applies the
	Wedderburn--Artin decomposition
	(\cite[Equation~(1.39)]{Wolf2012Quantum})
	to read off the block structure directly.
	The block decomposition is obtained by iterating
	the invariant-projection splitting and using strong
	induction on~$D$ to ensure termination.
\end{proof}

\begin{theorem}[Unitary diagonal positive fixed point in TP gauge]\label{thm:form_II}
	\lean{MPSTensor.exists_unitary_diag_posDef_fixedPoint_of_TP_of_isIrreducibleTensor}
	\leanok
	\uses{def:irreducible_tensor, def:tp_kraus}
	Let $A$ be an irreducible MPS tensor satisfying
	$\sum_i (A^i)^\dagger A^i = \Id$, and assume $D > 0$.
	Then there exist a unitary $U$ and a diagonal positive
	definite matrix $\Lambda$ such that, for $B^i := U^\dagger A^i U$,
	\[
		\sum_i (B^i)^\dagger B^i = \Id
		\qquad\text{and}\qquad
		\E_B(\Lambda) = \Lambda.
	\]
\end{theorem}

\begin{proof}
	\leanok
	\uses{thm:psd_fp_pd_irred}
	The TP hypothesis makes $\E_A$ a channel, so it has a nonzero
	PSD fixed point.
	The irreducible tensor condition implies that $\E_A$ is
	irreducible as a CP map, and Theorem~\ref{thm:psd_fp_pd_irred}
	upgrades the fixed point to a positive definite one.
	The spectral theorem diagonalizes this matrix by a unitary
	conjugation, and unitary conjugation preserves the TP
	normalization.
\end{proof}

The Burnside-style passage from irreducible tensors to full algebra spanning is
developed alongside the Wielandt material in Chapter~\ref{ch:wielandt};
see in particular Theorem~\ref{thm:irreducible_action_cumulative_span}.

\section{Canonical form from primitivity}\label{sec:canonical_from_primitivity}

The canonical form conditions
(Definition~\ref{def:cf_paper})
include an explicit hypothesis
that the self-overlap $O_{A_k A_k}(N) \to 1$
for each block.
The single source-faithful result of this section is
Theorem~\ref{thm:blocking_primitivity}: blocking an irreducible TP block by the
least common multiple of its peripheral periods makes the resulting transfer map
primitive, removing periodicity.

\begin{remark}[Two routes to the canonical-form predicate from normal blocks]\leanok
	Given a family of injective, left-canonical, strictly-ordered
	nonzero-weight blocks $(\mu_k, A_k)$, the canonical-form conditions of
	Definition~\ref{def:cf_paper} can be discharged from either of two
	spectral hypotheses on the per-block transfer map $\E_{A_k}$: a
	positive-definite primitive fixed point, or the peripheral-spectrum
	condition $\sigma_\partial(\E_{A_k}) = \{1\}$.  In both cases the
	overlap clause $O_{A_k A_k}(N) \to 1$ follows from the spectral-gap
	criterion of Chapter~\ref{ch:spectral}.  Both statements assume the
	blocks are already in the required normal form; they are not the
	arbitrary-input canonical-form existence theorem of
	\cite[Theorem~\texttt{Th:TIcanonical}]{PerezGarcia2007Matrix}.
\end{remark}

\begin{remark}\label{rem:primitivity_reduces_assumptions}\leanok
	The remark above shows that the overlap convergence hypothesis in the
	canonical form conditions is redundant once peripheral-spectrum
	primitivity is known.
	The argument factors through the spectral-gap formulation of
	Chapter~\ref{ch:spectral}: peripheral-spectrum primitivity is first
	converted to that spectral-gap condition, and the overlap limit is then
	deduced from it.
	In \cite{Cirac2021Matrix}, primitivity is obtained from
	irreducibility plus periodicity removal by blocking to period one.
	Theorem~\ref{thm:blocking_primitivity} is precisely this
	periodicity-removal step under irreducibility and TP normalization.
\end{remark}

\begin{theorem}[Blocking yields a peripheral-spectrum primitive transfer map]\label{thm:blocking_primitivity}
	\lean{MPSTensor.exists_blockTensor_isPrimitive_of_TP_of_isIrreducibleTensor}
	\leanok
	\uses{def:blocked_tensor, def:transfer_map, def:primitive_channel,
		def:tp_kraus, def:irreducible_tensor}
	Let $A$ be an irreducible MPS tensor with $D > 0$.
	Assume that $A$ satisfies the TP normalization
	$\sum_i (A^i)^\dagger A^i = \Id$.
	Then there exists a blocking length $p > 0$ such that the transfer map
	$\E_{A^{[p]}}$ of the blocked tensor $A^{[p]}$ is primitive,
	i.e. its peripheral spectrum is $\{1\}$.
\end{theorem}

\begin{proof}\leanok
	\uses{thm:peripheral_roots_irred_fp, thm:peripheral_finite,
		lem:common_power_eq_one, lem:peripheral_pow_singleton}
	The period-$p$ blocking step is the diagrammatic move
	of grouping the chain into consecutive blocks of length~$p$.
	Pass to the conjugate-transposed Kraus family $K_i := (A^i)^\dagger$.
	The TP normalization of $A$ makes $K$ unital, and irreducibility of $A$
	gives irreducibility of the associated Kraus map.
	The TP fixed-point theorem for $\E_A$ provides a positive definite fixed
	point for the adjoint map of $K$, so Theorem~\ref{thm:peripheral_roots_irred_fp}
	shows that every peripheral eigenvalue is a root of unity.
	Since $\mathrm{peripheral}(\E_A)$ is finite
	(Theorem~\ref{thm:peripheral_finite}),
	Lemma~\ref{lem:common_power_eq_one} gives a common exponent $p$
	such that $\mu^p = 1$ for all $\mu \in \mathrm{peripheral}(\E_A)$.
	Lemma~\ref{lem:peripheral_pow_singleton} then implies
	$\mathrm{peripheral}(\E_A^p) = \{1\}$.
	Finally, blocking corresponds to taking a power of the transfer map,
	so $\E_{A^{[p]}}$ is primitive.
\end{proof}

\begin{remark}\label{rem:blocking_primitivity_appendixA}\leanok
	Theorem~\ref{thm:blocking_primitivity} is the ``periodicity removal by blocking''
	step in \cite[Appendix~A]{Cirac2017MPDO_AnnPhys} for an irreducible block
	already placed in TP gauge. The later existence theorems then reuse
	this step.
\end{remark}

\section{Normal canonical form}\label{sec:normal_canonical_form}

\begin{definition}[Normal canonical form conditions]\label{def:is_normal_canonical_form}
		\lean{MPSTensor.IsNormalCanonicalForm}\leanok
		\uses{def:irreducible_tensor, def:primitive_channel, def:tp_kraus}
		A collection of scaling factors $(\mu_k)_{k=1}^r$ and block tensors $(A_k)_{k=1}^r$
		satisfies the \emph{normal canonical form conditions} if:
		\begin{enumerate}
			\item each $A_k$ is irreducible,
			\item each $A_k$ satisfies the TP normalization
			      $\sum_i (A_k^i)^\dagger A_k^i = \Id$,
			\item each transfer map $\E_{A_k}$ is primitive
			      (equivalently, has peripheral spectrum $\{1\}$),
			\item the moduli are non-increasing:
			      $|\mu_1| \ge |\mu_2| \ge \cdots \ge |\mu_r|$,
			\item all $\mu_k \neq 0$,
			\item all bond dimensions are positive.
		\end{enumerate}
		Here ``normal'' is used in the spectral sense of
		\cite{Cirac2017MPDO_AnnPhys}: irreducible blocks whose TP-normalized
		transfer maps are primitive. As noted in the remark after
		Definition~\ref{def:normal}, this is equivalent to the eventual
		block-injectivity formulation of Definition~\ref{def:normal}; see
		\cite[Proposition~3]{Sanz2010Quantum}.
\end{definition}

\begin{theorem}[Self-overlap is derived in normal canonical form]%
		\label{thm:ncf_overlap_tendsto_one}
		\lean{MPSTensor.IsNormalCanonicalForm.overlap_tendsto_one}\leanok
		\uses{def:is_normal_canonical_form, def:mpv_overlap}
		If $(\mu_k, A_k)$ satisfies the normal canonical form conditions, then
		$O_{A_k A_k}(N) \to 1$ for every block~$k$.
\end{theorem}

\begin{proof}\leanok
		\uses{def:is_normal_canonical_form}
		Peripheral-spectrum primitivity is part of the data. Together with
		irreducibility and TP normalization, it yields the spectral-gap
		primitive theorem needed for overlap convergence, so the self-overlap
		condition follows from the other hypotheses rather than appearing as
		a separate assumption.
\end{proof}

\begin{remark}
	See Theorem~\ref{thm:modulus_one_gauge_irr_tp} for the formalized modulus-one eigenvalue rigidity statement for irreducible trace-preserving blocks.
\end{remark}

\begin{lemma}[Mixed-transfer spectral radius from unit-modulus overlap]%
		\label{lem:overlap_norm_one_implies_spectral_radius_ge_one}
		\lean{MPSTensor.mixedTransferSpectralRadius_ge_one_of_mpvOverlap_norm_tendsto_one}
		\leanok
		\uses{def:mpv_overlap}
		Let $A$ and $B$ be MPS tensors of the same bond dimension.  If
		$\|\langle V^{(N)}(A)\,|\,V^{(N)}(B)\rangle\|\to 1$ as $N\to\infty$,
		then the mixed transfer map has spectral radius at least $1$.
\end{lemma}

\begin{proof}\leanok
		If the mixed transfer
		spectral radius were strictly less than~$1$, the iterated mixed transfer
		map would tend to zero in operator norm, hence so would its trace, and
		via the trace identity
		$\operatorname{tr}(F_{AB}^N) = \langle V^{(N)}(A), V^{(N)}(B)\rangle$
		the overlap itself would tend to~$0$, contradicting the hypothesis.
\end{proof}

\begin{theorem}[Gauge-phase equivalence from unit-modulus overlap]%
		\label{thm:nt_overlap_norm_one_gauge}
		\lean{MPSTensor.gaugePhaseEquiv_of_overlap_norm_tendsto_one_of_irreducible_TP}
		\leanok
		\uses{def:irreducible_tensor, def:tp_kraus, def:gauge_phase_equiv,
			def:mpv_overlap}
		Let $A$ and $B$ be irreducible MPS tensors of the same bond dimension
		satisfying $\sum_i (A^i)^\dagger A^i = \Id$ and
		$\sum_i (B^i)^\dagger B^i = \Id$.  If
		$\|\langle V^{(N)}(A)\,|\,V^{(N)}(B)\rangle\|\to 1$, then $A$ and~$B$
		are gauge-phase equivalent.
		% Same-bond-dimension case of \cite[Lemma~equalMPS,
		% lines~1080--1091]{Cirac2016MPDO_arXiv}; proportionality of MPVs is not assumed.
\end{theorem}

\begin{proof}\leanok
		\uses{lem:overlap_norm_one_implies_spectral_radius_ge_one,
			thm:modulus_one_gauge_irr_tp}
		By Lemma~\ref{lem:overlap_norm_one_implies_spectral_radius_ge_one},
		the mixed transfer spectral radius is at least~$1$.
		Theorem~\ref{thm:modulus_one_gauge_irr_tp} then gives gauge-phase
		equivalence.
\end{proof}

\section{Steps toward canonical form existence}%
\label{sec:canonical_pipeline_1606}

This section combines the components of the canonical-form
construction following \cite[Section~2.3 and Appendix~A]{Cirac2017MPDO_AnnPhys}.
The approach here requires blocking to remove non-trivial periodicity.
A separate reduction for periodic blocks without blocking
is developed in \cite{DeLasCuevas2017Irreducible}; see
Chapter~\ref{ch:periodic_ft_apps} for that direct periodic extension.

\begin{theorem}[Left-canonical normalization for irreducible TP blocks]%
	\label{thm:CFII_data}
	\lean{MPSTensor.exists_CFII_data_of_TP_of_isIrreducibleTensor}
	\leanok
	\uses{def:irreducible_tensor}
	Let $A$ be an irreducible MPS tensor already in TP
	gauge ($\sum_i (A^i)^\dagger A^i = \Id$) with $D > 0$.
	Then there exist a unitary~$U$ and a diagonal positive
	definite matrix~$\Lambda$ such that:
	\begin{enumerate}
		\item $B^i := U^\dagger A^i U$ is TP,
		\item $\E_B(\Lambda) = \Lambda$,
		\item $\Lambda$ is diagonal and positive definite,
		\item $A$ and $B$ generate the same MPV family.
	\end{enumerate}
	This is the left-canonical normalization step (Canonical Form~II)
	of \cite[Appendix~A]{Cirac2017MPDO_AnnPhys}. The unitary
	conjugation occurs only after TP normalization has already been
	achieved; it is a second step inside TP gauge, not a unitary
	normalization of an arbitrary irreducible block.
\end{theorem}

\begin{proof}\leanok
	\uses{thm:form_II, thm:samempv_unitary}
	Apply Theorem~\ref{thm:form_II} to obtain the unitary and
	the diagonal PD fixed point.
	Unitary conjugation preserves trace-preservation and the MPV
	family (Theorem~\ref{thm:samempv_unitary}).
\end{proof}

\begin{theorem}[Irreducible block decomposition with left-canonical normalization]%
	\label{thm:irreducible_block_decomp_with_cfii}
	\lean{MPSTensor.exists_irreducible_blockDecomp_with_CFII}
	\leanok
	\uses{def:irreducible_tensor}
	For any MPS tensor~$A$, there exist irreducible blocks
	$(A_k)_{k=1}^r$ that generate the same MPV family as~$A$.
	Moreover, for each block~$k$,
	if that block has already been placed in TP gauge and $D_k > 0$,
	one obtains a left-canonical normalization (unitary, diagonal PD fixed point,
	MPV preservation).
	This yields the left-canonical normalization; the further reduction to
	normal canonical form is carried out in the subsequent spectral analysis.
\end{theorem}

\begin{proof}\leanok
	\uses{thm:irred_decomp, thm:CFII_data}
	Combine Theorems~\ref{thm:irred_decomp}
	and~\ref{thm:CFII_data}.
\end{proof}

\begin{theorem}[Normal canonical form from a primitive block decomposition]%
		\label{thm:exists_normal_canonical_form}
		\lean{MPSTensor.exists_normalCanonicalForm_of_primitive_blockDecomp}
		\leanok
		\uses{def:is_normal_canonical_form, def:blocked_tensor}
		Suppose $A$ generates the same MPV family as a weighted block-diagonal
		tensor $\bigoplus_k \mu_k A_k$, and assume that:
		\begin{enumerate}
			\item each $A_k$ is irreducible,
			\item each $A_k$ satisfies the TP normalization
			      $\sum_i (A_k^i)^\dagger A_k^i = \Id$,
			\item each transfer map $\E_{A_k}$ is primitive
			      (equivalently, has peripheral spectrum $\{1\}$),
			\item the moduli $|\mu_k|$ are pairwise distinct,
			\item all $\mu_k \neq 0$,
			\item all bond dimensions are positive.
		\end{enumerate}
		Then there exist a blocking length $p > 0$, weights $(\nu_k)$, and blocked
		tensors $(B_k)$ such that $A^{[p]}$ generates the same MPV family as
		$\bigoplus_k \nu_k B_k$, and $(\nu_k, B_k)$ satisfies the normal canonical
		form conditions.
\end{theorem}
% Scope restriction (prepared primitive block decomposition): not the
% full source theorem (lines 742--763), which starts from
% an arbitrary translation-invariant MPS representation. See
% docs/paper-gaps/pgvwc07_ti_canonical_form_scope.tex.

\begin{proof}\leanok
		\uses{def:is_normal_canonical_form}
		No new blocking is needed here: the given blocks are already primitive,
		so one may take $p = 1$.
		The theorem is stated with a blocking length because earlier reduction
		steps may require a common blocking period; under these hypotheses
		that period can be chosen to be $1$.
		Since the moduli $|\mu_k|$ are pairwise distinct, there is a reordering
		of the blocks for which they become strictly decreasing.
		After this reordering, set $\nu_k := \mu_k$ and $B_k := A_k$ in the new
		order.
		Irreducibility, TP normalization, primitivity, nonzero weights, and
		positive bond dimensions are unchanged by the reordering, so the
		reordered family satisfies Definition~\ref{def:is_normal_canonical_form}
		and generates the same MPV family as~$A$.
\end{proof}

\begin{remark}[Scope of Theorem~\ref{thm:exists_normal_canonical_form}]
  Theorem~\ref{thm:exists_normal_canonical_form} assumes all six block
  hypotheses---irreducibility, TP normalization, primitivity, distinct
  nonzero weight moduli, and positive bond dimensions---in advance, and
  produces the normal canonical form by reordering.  It is \emph{not} the
  arbitrary-input canonical-form existence theorem of
  \cite[lines~237--250]{Cirac2017MPDO_AnnPhys}, which starts from an
  arbitrary tensor and derives the primitive block structure.  That
  arbitrary-input direction is the open formalization gap documented in
  Section~\ref{subsec:paperbnt_supplier}.
\end{remark}

\section{TP-gauge reduction for arbitrary tensors}%
\label{sec:tp_gauge_arbitrary}

\begin{theorem}[TP-gauge reduction for arbitrary tensors with zero-block separation]%
	\label{thm:tp_gauge_arbitrary}
	\lean{MPSTensor.exists_tp_gauge_from_arbitrary_with_zeroTail}
	\leanok
	\uses{def:irreducible_tensor, def:zero_mps_tensor}
	For any MPS tensor $A$, there exist:
	\begin{enumerate}
		\item a zero-tail bond dimension $D_0$,
		\item nontrivial blocks $(B_k)_{k=1}^r$ with nonzero weights
		      $(\mu_k)_{k=1}^r$, each block irreducible and
		      left-canonical ($\sum_i (B_k^i)^\dagger B_k^i = \Id$),
		      with positive bond dimension,
	\end{enumerate}
	such that
	\[
		V^{(N)}(A)_\sigma
		= V^{(N)}(\mathbf{0}_{d,D_0})_\sigma
		+ V^{(N)}\!\bigl(\textstyle\bigoplus_{k=1}^r \mu_k B_k\bigr)_\sigma.
	\]
	This is the furthest unconditional reduction step available for arbitrary tensors
	before periodicity removal.
\end{theorem}

\begin{proof}\leanok
	\uses{thm:zero_block_separation, thm:tp_data_irreducible}
	Apply Theorem~\ref{thm:zero_block_separation} to separate the
	trivial block and nontrivial blocks. Then apply the blockwise TP-gauge theorem
	to the nontrivial blocks. The MPV identity chains through both steps.
\end{proof}

\section{Blocking and primitivity reduction}%
\label{sec:blocking_infrastructure}

The intermediate constructions that transport MPV equalities, weights, transfer
maps, and physical-alphabet casts through blocking by an integer~$p$ are
collected in the Lean modules \texttt{MPS/Core/BlockingInfrastructure.lean} and
\texttt{MPS/Core/BlockingTransfer.lean}.  The single source-facing consequence
needed downstream is the existence of a common blocking period under which
every transfer map in a finite family becomes primitive.

\begin{theorem}[Common blocking period for a block family]%
	\label{thm:common_blocking_period}
	\lean{MPSTensor.exists_common_blocking_all_primitive}
	\leanok
	\uses{def:blocked_tensor}
	Let $(B_k)_{k=1}^r$ be a finite family of MPS tensors with
	positive bond dimensions, each admitting a blocking period~$p_k$
	making its transfer map primitive.
	Then $p = \operatorname{lcm}(p_1, \ldots, p_r)$ is a common
	period: $\E_{B_k^{[p]}}$ is primitive for every~$k$.
\end{theorem}

\begin{proof}\leanok
	Each $p_k$ divides the least common multiple~$p$, so primitivity of $\E_{B_k^{[p_k]}}$
	transports to primitivity of $\E_{B_k^{[p]}}$ via the standing divisibility
	identity for transfer-map powers.
\end{proof}

\section{Period removal (cyclic-sector decomposition)}%
\label{sec:cyclic_sectors}

In the non-periodic route of \cite{Cirac2017MPDO_AnnPhys} the only required step
is periodicity removal by blocking: after blocking each irreducible TP block by
the least common multiple of its peripheral periods, the resulting tensor has no
nontrivial $p$-periodic vectors
\cite[Section~2.3]{Cirac2017MPDO_AnnPhys}.  The cyclic-sector decomposition
itself comes from the peripheral spectrum of the adjoint transfer map
\cite[Theorem~6.6]{Wolf2012Quantum}.  Three source-facing statements summarize
this step: existence of the cyclic decomposition after blocking, the version
specialized to irreducible TP tensors, and the strengthened version where the
resulting sector blocks are themselves primitive and irreducible.

\begin{theorem}[Period removal by blocking (1606.00608, §2.3)]%
	\label{thm:cyclic_sector_decomp_after_blocking}
	\lean{MPSTensor.exists_cyclic_sector_decomp_after_blocking}
	\leanok
	\uses{def:blocked_tensor}
	Let $A$ be a left-canonical irreducible tensor of period $m$.
	Then the blocked tensor $A^{[m]}$ admits:
	\begin{enumerate}
		\item left-canonical sector tensors $(C_k)_{k=1}^m$,
		\item orthogonal projections $(P_k)_{k=1}^m$ summing to $\Id$ and satisfying
		      $\E_{A^\dagger}(P_{k+1}) = P_k$,
		\item compression linear equivalences
		      $\varphi_k : \MN{\dim C_k} \xrightarrow{\sim} P_k \MN{D} P_k$,
		\item a direct-sum MPV decomposition
		      $A^{[m]} \sim \bigoplus_{k=1}^m C_k$.
	\end{enumerate}
	This is the period-removal step: after blocking by $m$, the
	result admits a unit-weight decomposition with no non-trivial
	$p$-periodic vectors.  The cyclic projections $P_k$ come from
	the peripheral spectrum of the adjoint transfer map
	(\cite[Theorem~6.6]{Wolf2012Quantum}).
\end{theorem}

\begin{proof}\leanok
	The cyclic decomposition of the adjoint transfer map gives projections
	$(P_k)$ with $\E_{A^\dagger}(P_{k+1}) = P_k$.  Iterating the cyclic relation $m$
	times shows $(\E_{A^\dagger})^m(P_k) = P_k$, and the blocked-transfer identity
	$\E_{(A^{[m]})^\dagger} = (\E_{A^\dagger})^m$ identifies these as the adjoint
	transfer map of the blocked tensor.  Apply the block decomposition from
	adjoint-fixed projections (formalized in
	\texttt{MPS/CanonicalForm/CyclicSectors/FixedAdjoint.lean}).
\end{proof}

\begin{theorem}[Period removal for irreducible TP tensors]%
	\label{thm:cyclic_sector_decomp_irr_tp}
	\lean{MPSTensor.exists_cyclic_sector_decomp_of_TP_of_isIrreducibleTensor}
	\leanok
	\uses{def:blocked_tensor, def:irreducible_tensor}
	Let $A$ be a left-canonical irreducible tensor.
	Then there exist $m \ge 1$ and left-canonical blocks
	$(C_k)_{k=1}^m$ such that the $m$-blocked tensor $A^{[m]}$
	admits a unit-weight decomposition into the $C_k$.
	This is the period-removal step of \cite[Section~2.3]{Cirac2017MPDO_AnnPhys}.
\end{theorem}

\begin{proof}\leanok
	\uses{thm:cyclic_sector_decomp_after_blocking}
	Left-canonicality gives $\sum_i (A^i)^\dagger A^i = \Id$, so the
	conjugate-transposed family $K_i := (A^i)^\dagger$ is unital; irreducibility
	passes to the adjoint map, and the quantum Perron--Frobenius theorem provides
	a positive-definite fixed point.  The cyclic peripheral-spectrum theorem then
	provides a period~$m$ and the cyclic spectral data, and
	Theorem~\ref{thm:cyclic_sector_decomp_after_blocking} produces the
	decomposition.
\end{proof}

\begin{theorem}[Period removal with primitive irreducible sectors]%
\label{thm:cyclic_sector_decomp_irr_tp_prim_irr}
	\lean{MPSTensor.exists_primitive_irreducible_cyclic_sector_decomp_of_TP_of_isIrreducibleTensor}
	\leanok
	\uses{thm:cyclic_sector_decomp_after_blocking, thm:blocked_sector_unconditional}
	Let $A$ be a left-canonical irreducible tensor. Then there is a period
	$m \ge 1$ and a unit-weight decomposition of $A^{[m]}$ into sector blocks
	$(C_k)_{k=1}^m$ such that every $C_k$ is left-canonical, has primitive
	transfer map, is tensor-irreducible, and has nonzero bond dimension.

	Equivalently: the cyclic-sector decomposition of $A^{[m]}$ produces unit
	blocks that already satisfy the trace-preservation, primitivity, and
	irreducibility hypotheses needed for the common-blocking step
	(Lemma~\ref{thm:exists_common_blocked_cyclic_sector_family}).
\end{theorem}

\begin{proof}\leanok
	\uses{thm:peripheral_cyclic_structure,
		thm:cyclic_sector_decomp_after_blocking, thm:blocked_sector_unconditional}
	The proof repeats the cyclic-sector construction of
	Theorem~\ref{thm:cyclic_sector_decomp_irr_tp}, keeps the projection and
	compression data, and applies
	Theorem~\ref{thm:blocked_sector_unconditional} to the resulting
	sector blocks.
\end{proof}

\section{Reduction to primitive blocks}%
\label{sec:reduction_to_normal_form}

The structural after-blocking statements that take an arbitrary tensor or a
TP-primitive-irreducible block to a primitive block decomposition appear in
Chapter~\ref{ch:canonical_after_blocking}, and the block-injectivity statement
in Chapter~\ref{ch:wielandt}.  The single statement recorded in this chapter
is the normality of a TP-primitive irreducible block.

\begin{theorem}[TP-primitive irreducible tensor is normal]%
	\label{thm:normal_tp_primitive_irred}
	\lean{MPSTensor.isNormal_of_tp_primitive_irreducible}
	\leanok
	\uses{def:primitive_mps, def:irreducible_tensor, def:normal}
	Let $A$ be a left-canonical MPS tensor with $D > 0$,
	irreducible transfer map, and peripheral spectrum $\{1\}$.
	Then $A$ is normal.
\end{theorem}

\begin{proof}\leanok
	\uses{thm:normal_from_primitive_posdef, thm:psd_fp_pd_irred}
	Peripheral primitivity and irreducibility yield the spectral-gap
	condition.  The PSD fixed point is positive definite by
	Theorem~\ref{thm:psd_fp_pd_irred}.
	Apply Theorem~\ref{thm:normal_from_primitive_posdef}.
\end{proof}

\subsection{Basis of Normal Tensors}

\begin{theorem}[Sorting by strictly decreasing weight norms]
	\label{thm:bnt_sorted_reindexing}
	Let $(\mu_k)_{k=1}^r$ be nonzero complex weights whose moduli are pairwise
	distinct.  Then there exists a permutation $\sigma$ of $\{1,\dots,r\}$ such that
	\[
		|\mu_{\sigma(1)}| > |\mu_{\sigma(2)}| > \cdots > |\mu_{\sigma(r)}| > 0.
	\]
\end{theorem}

\begin{theorem}[Sorted decomposition preserves the represented MPS family]
	\label{thm:bnt_sorted_blockdecomp}
	\uses{thm:bnt_sorted_reindexing}
	Under the same hypothesis, with $\sigma$ as above and
	$\tilde\mu_k := \mu_{\sigma(k)}$, $\tilde B_k := B_{\sigma(k)}$,
	\[
		\mc{V}\!\bigl(\textstyle\bigoplus_k \mu_k B_k\bigr)
		\;=\; \mc{V}\!\bigl(\textstyle\bigoplus_k \tilde\mu_k \tilde B_k\bigr),
		\qquad |\tilde\mu_1| > \cdots > |\tilde\mu_r|.
	\]
\end{theorem}

\begin{theorem}[Normal canonical form after sorting]
	\label{thm:bnt_sorted_ncf}
	\uses{def:is_normal_canonical_form, thm:bnt_sorted_blockdecomp}
	Suppose $(\mu_k, B_k)_{k=1}^r$ satisfies $\mu_k \neq 0$, $|\mu_k|$ pairwise
	distinct, and each $B_k$ is irreducible, left-canonical
	($\sum_i (B_k^i)^\dagger B_k^i = \Id$), with primitive transfer map and
	$\dim B_k > 0$.  Then after sorting by strictly decreasing
	$|\mu_{\sigma(k)}|$, the resulting family $(\tilde\mu_k, \tilde B_k)$
	satisfies the normal canonical form conditions of
	Definition~\ref{def:is_normal_canonical_form}.
\end{theorem}

\begin{definition}[Granular sector decomposition]
	\label{def:trivial_sector_decomp}
	\lean{MPSTensor.trivialSectorDecomp}
	\leanok
	\uses{def:paper_sector_data}
	Given nonzero weights $(\mu_k)_{k=1}^r$ and blocks $(B_k)_{k=1}^r$, the
	\emph{granular sector decomposition} has $r$ sectors, one basis block
	$B_k$ per sector, multiplicity~$1$ in each sector, and sector weight
	$\mu_k$ on the $k$th sector.  Equivalently, as a finite sum,
	\[
		\sum_{k=1}^{r} \mu_k^{\,N} V^{(N)}(B_k)_\sigma
		\;=\; V^{(N)}\!\bigl(\textstyle\bigoplus_k \mu_k B_k\bigr)_\sigma.
	\]
\end{definition}

\begin{lemma}[Granular sector decomposition preserves the represented family]
	\label{lem:trivial_sector_decomp_same_mpv}
	\lean{MPSTensor.sameMPV₂_trivialSectorDecomp}
	\leanok
	\uses{def:trivial_sector_decomp}
	The granular sector decomposition associated to $(\mu_k,B_k)_k$ represents
	the same MPS family as the original weighted block sum
	$\bigoplus_k \mu_k B_k$.
\end{lemma}

\begin{proof}\leanok
	\uses{def:trivial_sector_decomp}
	Expanding the sector decomposition gives the sum
	$\sum_k \mu_k^N V^{(N)}(B_k)_\sigma$, which is exactly the finite
	block-sum expression for $\bigoplus_k \mu_k B_k$.
\end{proof}

\begin{theorem}[Trivial sector decomposition for the sorted distinct-norm case]
	\label{thm:bnt_trivial_sector}
	\uses{def:trivial_sector_decomp, lem:trivial_sector_decomp_same_mpv}
	If $\mu_k \neq 0$ for all $k$ and $|\mu_1| > \cdots > |\mu_r|$, then the
	granular sector decomposition $P = \mathrm{trivialSectorDecomp}(\mu,B)$
	satisfies $V^{(N)}(P) = V^{(N)}(\bigoplus_k \mu_k B_k)$ for every
	$N \ge 1$ and has exactly one copy per sector.
\end{theorem}

\begin{definition}[Norm-class partition]
	\label{def:norm_class_grouping_data}
	Let $(\mu_k)_{k=1}^r$ be a finite family of complex weights.  A
	\emph{norm-class partition} consists of: a number of classes $s \le r$;
	representative values $\lambda_1 > \cdots > \lambda_s > 0$;
	multiplicities $(m_j)_{j=1}^s$ with $\sum_j m_j = r$; for each $j$ an
	enumeration $\iota_j : \{1,\dots,m_j\} \hookrightarrow \{1,\dots,r\}$ of
	the indices $k$ with $|\mu_k| = \lambda_j$; and the regrouping identity
	\[
		\sum_{k=1}^{r} f(k) \;=\; \sum_{j=1}^{s} \sum_{q=1}^{m_j} f(\iota_j(q))
	\]
	for every function $f$.
\end{definition}

\begin{definition}[Norm-class enumeration in decreasing order]
	\label{def:norm_class_grouping}
	\uses{def:norm_class_grouping_data}
	For any finite family $(\mu_k)_{k=1}^r$ of complex weights, sorting
	the distinct moduli $\{|\mu_k| : 1 \le k \le r,\ \mu_k \neq 0\}$ in
	strictly decreasing order yields a norm-class partition (in the
	sense of Definition~\ref{def:norm_class_grouping_data}) with
	representative values $\lambda_1 > \cdots > \lambda_s$.
\end{definition}

\section{Open directions}\label{sec:open_directions}

\begin{remark}[Canonical-form reduction before the equal-case comparison]%
\label{rem:canonical_pipeline_gaps}
	\leanok
	The canonical-form proof in the MPS literature first separates the nilpotent
	part (which contributes only to the length-zero trace) from the nonzero part,
	decomposes the nonzero part into irreducible blocks, places each block in TP
	gauge, blocks each period-$m$ block over $m$ sites, and restricts to its
	cyclic spectral sectors.  The sector tensors are trace-preserving, primitive,
	irreducible, and have positive bond dimension.  For a finite nonzero block
	family one chooses a common multiple of all periods so that every sector is
	expressed over the same blocked physical alphabet
	\cite[Section~2.3 and Appendix~A]{Cirac2017MPDO_AnnPhys}.

	The common-alphabet cyclic-sector family and the after-blocking
	zero-tail and common-sector transport statements are stated and proved
	in Chapter~\ref{ch:canonical_after_blocking}; the present chapter
	records only
	Theorems~\ref{thm:cyclic_sector_decomp_after_blocking},
	\ref{thm:cyclic_sector_decomp_irr_tp}, and
	\ref{thm:cyclic_sector_decomp_irr_tp_prim_irr}.
	The sorted normal-canonical-form theorem
	(Theorem~\ref{thm:bnt_sorted_ncf}) assumes pairwise distinct nonzero weight
	norms; equal-norm families require the basis-of-normal-tensors and
	multiplicity comparison developed in Chapter~\ref{ch:bnt}.
\end{remark}

\subsection{Material from PGVWC06 not covered by this chapter}%
\label{subsec:pgvwc06_out_of_scope}

This chapter follows the presentation of~\cite{Cirac2017MPDO_AnnPhys}
and~\cite{Cirac2021Matrix}, which works with site-independent
translation-invariant tensors, transfer-map primitivity, and CPM
fixed-point normalization.  The 2007 preprint~\cite{PerezGarcia2007Matrix}
contains eight named theorems in its canonical-form section~\S III
(lines~399--1177 of \texttt{MPSarchive.tex}).  Theorem~\texttt{Th:TIcanonical}
has a partial, conditional counterpart here; the opening paragraphs of this
chapter make the scope limitation explicit.  The remaining seven theorems
of \cite{PerezGarcia2007Matrix} \S III are not covered:

\begin{itemize}
  \item \textbf{Thm.~\texttt{OBC-Vidal} (lines~431--443): completeness and
    OBC canonical form.}
    Any pure state admits an OBC-MPS representation of bond dimension
    $D\le d^{\lfloor N/2\rfloor}$ with simultaneous left-normalization and
    a positive full-rank diagonal interleaving matrix~$\Lambda^{[m]}$.
    \emph{Not covered:} the chapter starts from a site-independent TI
    tensor; constructing the finite-OBC Schmidt canonical form requires a
    site-dependent Chain layer not present in this development.

  \item \textbf{Thm.~\texttt{free-OBC} (lines~466--486): OBC representation
    freedom.}
    Every OBC-MPS $\{B^{[m]}_i\}$ is related to the canonical one by
    site-local invertible gauge matrices $Y_j,Z_j$ with $Y_j Z_j = \mathbf{1}$.
    \emph{Not covered:} same reason as \texttt{OBC-Vidal}; local-gauge
    analysis for site-dependent open-boundary matrices is outside the TI-MPS
    scope.

  \item \textbf{Thm.\ ``Site-independent matrices'' (lines~620--630).}
    Every TI pure state with PBC on a finite ring of $N$ sites has a
    site-independent MPS representation; a bond-$D$ OBC representation
    lifts to one at bond dimension at most~$ND$.
    \emph{Not covered:} the chapter assumes a site-independent MPV tensor
    as input; the block-circulant construction producing it is not covered.

  \item \textbf{Thm.~\texttt{Th:periodic} (lines~849--858): periodic
    decomposition.}
    If the transfer map of a single-block TI canonical MPS has $p$
    eigenvalues of modulus~$1$, the state is a superposition of $p$
    $p$-periodic states of bond dimension~$D$, and vanishes when $p \nmid N$.
    \emph{Partially formalized:} the cyclic-sector machinery of
    Section~\ref{sec:cyclic_sectors} (Theorem~\ref{thm:cyclic_sector_decomp_after_blocking})
    removes the period and yields primitive sector tensors, but the explicit
    state-vector decomposition into $p$-periodic summands is not separately
    stated.

  \item \textbf{Thm.\ ``Interpretation of $\Lambda$'' (lines~987--993).}
    Under condition~C2, the eigenvalues of the half-chain reduced density
    operator converge as $N\to\infty$ to $\Lambda\otimes\Lambda$ with
    $\Lambda$ from Theorem~\texttt{Th:TIcanonical}.
    \emph{Not covered:} Theorem~\ref{thm:CFII_data} produces $\Lambda$, but the
    asymptotic eigenvalue statement requires correlation-length and entropy
    material not addressed here.

  \item \textbf{Uniqueness of the canonical form
    (\cite[lines~1002--1015]{PerezGarcia2007Matrix}).}
    Under condition~C1 with unique OBC canonical form, two TI canonical
    $D$-MPS representations of the same state satisfy $B_i = U C_i U^\dagger$
    for a global unitary $U$.
    \emph{Not covered in this chapter.} A modern analogue is proved in
    Chapter~\ref{ch:ft_proof} (global-gauge equality theorem); the original
    finite-ring uniqueness of \cite{PerezGarcia2007Matrix} under C1/C2 and
    $N > 2L_0 + D^4$ is not covered here.

  \item \textbf{Thm.\ ``Obtaining the TI canonical form''
    (lines~1154--1165).}
    The site-independent $D\times D$ matrices can be recovered from an OBC
    canonical form by solving an $N$-independent system~(S) of $O(D^8)$
    quadratic equations; every solution is related to the canonical form by
    a global unitary.
    \emph{Not covered:} this chapter addresses existence and uniqueness
    abstractly; numerical and algorithmic recovery is not formalized.
\end{itemize}

\subsection{Building a BNT from normal blocks with normalized weights}%
\label{subsec:paperbnt_supplier}

The BNT predicate (Definition~\ref{def:bnt_paper}, formalized as
\texttt{IsBNTCanonicalForm} in Chapter~\ref{ch:bnt}) bundles the
basis-of-normal-tensors structure of \cite{Cirac2017MPDO_AnnPhys} with
multiplicity and span fields.  The following theorem starts from a family
of normalized block data and produces a sector decomposition satisfying that
predicate.

\begin{theorem}[BNT from normal blocks with normalized weights]%
  \label{thm:paperbnt_supplier_prepared}
  \lean{MPSTensor.exists_isBNTCanonicalForm_of_tp_primitive_irr_injective_blocks}
  \leanok
  Let $r \in \mathbb{N}$, let $\dim : \mathrm{Fin}\,r \to \mathbb{N}$,
  let $\mu : \mathrm{Fin}\,r \to \mathbb{C}$, and let
  $\mathrm{blocks}_k : \texttt{MPSTensor}\,d\,(\dim k)$ for $k < r$.
  Assume:
  \begin{enumerate}
    \item \textbf{positive bond dimensions:} $0 < \dim k$ for all $k$,
    \item \textbf{TP normalization:} each block satisfies
          $\sum_i (\mathrm{blocks}_k^i)^\dagger\,\mathrm{blocks}_k^i = \Id$
          (\texttt{IsLeftCanonical}),
    \item \textbf{primitivity:} the transfer map $\mathcal{E}_{\mathrm{blocks}_k}$
          is primitive for each $k$,
    \item \textbf{irreducibility:} each block satisfies
          \texttt{IsIrreducibleTensor},
    \item \textbf{injectivity:} each block satisfies \texttt{IsInjective},
    \item \textbf{nonzero weights:} $\mu_k \neq 0$ for all $k$,
    \item \textbf{bounded weights:} $|\mu_k| \le 1$ for all $k$,
    \item \textbf{global unit witness:} $\exists\,k,\;|\mu_k| = 1$.
  \end{enumerate}
  Then there exists a \texttt{SectorDecomposition} $P$ such that
  \begin{itemize}
    \item $P.\texttt{toTensor}$ and
          $\texttt{toTensorFromBlocks}\,\mu\,\mathrm{blocks}$
          generate the same MPV family at every length
          (\texttt{SameMPV\textsubscript{2}}), and
    \item $P$ satisfies \texttt{IsBNTCanonicalForm}.
  \end{itemize}
\end{theorem}

\begin{proof}\leanok
  The phase-class quotient collapses gauge-phase-equivalent blocks into
  sectors with unit-modulus per-copy scalars (derived from the TP + primitive +
  irreducible hypotheses via the MPV phase-class infrastructure).
  The resulting sector decomposition inherits the BNT data fields from
  \texttt{exists\_bnt\_sectorDecomp\_of\_tp\_primitive\_irr\_blocks};
  the bounded-weight and global-unit hypotheses close the weight-norm fields
  of \texttt{IsBNTCanonicalForm} directly.
\end{proof}

\begin{remark}[Arbitrary-input direction: open formalization gap]%
  \label{rem:arbitrary_input_bridge_gap}
  Theorem~\ref{thm:paperbnt_supplier_prepared} assumes eight pre-normalized
  hypotheses on the block data; it does not produce a sector decomposition
  starting from an arbitrary MPS tensor.  Three independent obstacles
  separate it from the arbitrary-input statement.

  \begin{enumerate}
    \item \textbf{Zero-tail at $N=0$.}  The same-MPV equivalence used in the
      BNT predicate quantifies over all $N \ge 0$, so the $N=0$ clause
      forces the bond dimension of the sector decomposition to equal the
      bond dimension of the input.  The arbitrary-input reduction naturally
      yields equality only for $N \ge 1$ together with a separate
      length-zero dimension identity; matching the zero-tail bond dimension
      fails once all-zero blocks are removed by
      Theorem~\ref{thm:zero_block_separation}.

    \item \textbf{Global weight normalization.}  The BNT predicate
      requires $|\mu_k| \le 1$ for all $k$ and $|\mu_k| = 1$ for some $k$.
      These are normalization conventions, not consequences of an
      arbitrary tensor.  If a left-canonical block $C$ has unit weight,
      the scaled input $A = 2 \cdot C$ has weight of modulus $2$, violating
      the bound.  The hypotheses must therefore be supplied externally,
      or the statement weakened to a rescaled form.

    \item \textbf{Per-block unit-modulus witness.}  The BNT predicate
      requires
      $\forall j,\;\exists q,\;\bigl|\mathrm{weight}\,j\,q\bigr| = 1$.
      For a decaying sector such as $C \oplus \tfrac{1}{2} D$, this fails:
      no weight of the $D$ sector has unit modulus.  The condition is
      therefore an explicit hypothesis rather than a consequence of the
      block decomposition.
  \end{enumerate}

  Closing the gap for arbitrary tensors remains open.
\end{remark}

\subsection{Auxiliary lemmas}%
\label{subsec:internal_lean_only_labels}

The lemmas below collect the small structural facts about transfer maps,
blocking, and cyclic-sector reblocking that are used in the after-blocking
constructions of Chapter~\ref{ch:ft_proof} and
Chapter~\ref{ch:canonical_after_blocking} and in the periodic extension of
Chapter~\ref{ch:periodic_ft_apps}.

\begin{lemma}[Canonical form from peripheral primitivity]%
\label{thm:canonical_from_primitive}
\uses{def:is_canonical_form, def:irreducible_tensor, def:tp_kraus}
Let $(\mu_k,A_k)_{k=1}^r$ be a family of nonzero weights and left-canonical
irreducible blocks with positive bond dimensions, with strictly decreasing
moduli $|\mu_1|>\cdots>|\mu_r|>0$ and $|\mu_1|=1$, and assume each transfer
map satisfies $\sigma_\partial(\E_{A_k}) = \{1\}$ (peripheral-spectrum
primitivity). Then $(\mu_k,A_k)_{k=1}^r$ satisfies the canonical-form
conditions of Definition~\ref{def:is_canonical_form}. In particular,
$O_{A_k A_k}(N) \to 1$ for every $k$.
\end{lemma}

\begin{lemma}[Adjoint primitivity equivalence]%
\label{thm:primitive_adjoint_equiv}
\lean{MPSTensor.IsPrimitive.adjoint_iff}\leanok
\uses{def:primitive_channel}
Let $\E$ be a finite-dimensional completely positive map. Then $\E$ is
primitive if and only if its Frobenius adjoint $\E^\dagger$ is primitive.
Equivalently, $\sigma_\partial(\E) = \{1\}$ iff
$\sigma_\partial(\E^\dagger) = \{1\}$, since the peripheral spectrum of
$\E^\dagger$ consists of the complex conjugates of the peripheral spectrum
of $\E$.
\end{lemma}

\begin{lemma}[Primitivity under blocking by a multiple]%
\label{thm:primitive_blocking_divisible}
\lean{MPSTensor.isPrimitive_transferMap_blockTensor_of_dvd}\leanok
\uses{def:blocked_tensor, def:primitive_channel}
Let $A$ be an MPS tensor with positive bond dimension and let $p_0,p\ge 1$
with $p_0 \mid p$. If $\E_{A^{[p_0]}}$ is primitive, then $\E_{A^{[p]}}$ is
primitive.
\end{lemma}

\begin{lemma}[Blocking the assembled weighted block sum]%
\label{thm:block_tensor_to_tensor_from_blocks}
\lean{MPSTensor.sameMPV₂_blockTensor_toTensorFromBlocks}\leanok
\uses{def:blocked_tensor, def:block_diagonal_tensor, def:same_mpv2}
For nonzero weights $(\mu_k)_{k=1}^r$, blocks $(B_k)_{k=1}^r$, and any
blocking length $p\ge 1$, the blocked assembled tensor agrees with the
assembled tensor of the blocked blocks with powered weights:
\[
  (\textstyle\bigoplus_k \mu_k B_k)^{[p]}
  \sim  \textstyle\bigoplus_k \mu_k^{\,p}\, B_k^{[p]},
\]
where $\sim$ denotes generation of the same MPV family at every length.
\end{lemma}

\begin{lemma}[Nonzero weights survive blocking]%
\label{lem:block_weights_ne_zero}
\lean{MPSTensor.blockWeights_ne_zero}\leanok
If $\mu \in \mathbb{C}$ is nonzero and $p \ge 1$, then $\mu^p \neq 0$. In
particular, blocking preserves the nonzero-weight hypothesis of a
weighted block decomposition.
\end{lemma}

\begin{lemma}[Primitive cyclic sectors of irreducible TP tensors]%
\label{thm:has_primitive_irreducible_cyclic_sectors_tp_irr}
\lean{MPSTensor.hasPrimitiveIrreducibleCyclicSectors_of_TP_of_isIrreducibleTensor}\leanok
\uses{def:irreducible_tensor, def:tp_kraus, def:primitive_channel,
      thm:cyclic_sector_decomp_irr_tp_prim_irr}
Let $A$ be a left-canonical irreducible MPS tensor with $D > 0$. Then there
exist a period $m \ge 1$ and a family $(C_k)_{k=1}^m$ of left-canonical
blocks with primitive transfer maps, tensor-irreducible structure, and
positive bond dimensions such that $A^{[m]}$ generates the same MPV family
as the unit-weight assembled tensor $\bigoplus_{k=1}^m C_k$.
\end{lemma}

\begin{lemma}[Common reblocking of cyclic-sector families]%
\label{thm:exists_common_blocked_cyclic_sector_family}
\lean{MPSTensor.exists_commonBlockedCyclicSectorFamily_of_hasPrimitiveIrreducibleCyclicSectors}\leanok
\uses{thm:has_primitive_irreducible_cyclic_sectors_tp_irr,
      def:blocked_tensor}
Given a finite family $(B_k)_{k=1}^r$ of left-canonical irreducible TP
tensors with positive bond dimensions and per-block period-removal lengths
$(m_k)_{k=1}^r$, there exists a common blocking length $p$ such that, for
every $k$, the blocked tensor $B_k^{[p]}$ admits a unit-weight assembled
decomposition over a single common blocked alphabet $d^{[p]}$ into
left-canonical blocks with primitive transfer maps, tensor irreducibility,
and positive bond dimensions.
\end{lemma}

\begin{lemma}[Nonzero unit weight of the common reblocked family]%
\label{thm:common_blocked_cyclic_sector_flat_weight_ne_zero}
\lean{MPSTensor.CommonBlockedCyclicSectorFamily.commonFlatWeight_ne_zero}\leanok
\uses{thm:exists_common_blocked_cyclic_sector_family}
The unit weight on each sector of the common reblocked cyclic-sector family
of Lemma~\ref{thm:exists_common_blocked_cyclic_sector_family} is nonzero.
\end{lemma}

\begin{lemma}[Sector irreducibility for the adjoint blocked map]%
\label{thm:sector_irred_unconditional}
\lean{MPSTensor.isIrreducibleOnCorner_of_cyclic_decomp_mps}\leanok
\uses{def:irreducible_tensor}
Let $A$ be a left-canonical irreducible MPS tensor with $D>0$, let $m$ be
the period of the cyclic-sector decomposition of $A$, and let
$(P_u)_{u=1}^m$ be the cyclic projections of the adjoint transfer map
satisfying $\E_{A^\dagger}(P_{u+1}) = P_u$ and $\sum_u P_u = \Id$. Then
the restriction of $(\E_{A^\dagger})^m$ to each corner $P_u \MN{D} P_u$ is
irreducible.
\end{lemma}

\begin{lemma}[Iterated blocking at the transfer-map level]%
\label{thm:iterated_blocking_transfer}
\lean{MPSTensor.transferMap_blockTensor_mul}\leanok
\uses{def:blocked_tensor, def:transfer_map}
For any MPS tensor $A$ and integers $m,n \ge 1$,
\[
  \E_{(A^{[m]})^{[n]}} \;=\; \E_{A^{[mn]}}.
\]
\end{lemma}

\begin{definition}[Primitive irreducible cyclic-sector data]%
\label{def:primitive_irreducible_cyclic_sectors}
\lean{MPSTensor.HasPrimitiveIrreducibleCyclicSectors}\leanok
\uses{def:blocked_tensor, def:irreducible_tensor, def:primitive_channel,
      def:tp_kraus}
$A$ \emph{has primitive irreducible cyclic sectors} of period $m \ge 1$ if
there exists a family $(C_k)_{k=1}^m$ of MPS tensors such that $A^{[m]}$
generates the same MPV family as $\bigoplus_{k=1}^{m} C_k$ (unit weights),
and each $C_k$ is left-canonical, has primitive transfer map, is
tensor-irreducible, and has positive bond dimension.
\end{definition}

\begin{lemma}[Common reblocking at a prescribed multiple]%
\label{thm:exists_common_blocked_cyclic_sector_family_common_multiple}
\lean{MPSTensor.exists_commonBlockedCyclicSectorFamily_of_commonMultiple}\leanok
\uses{thm:exists_common_blocked_cyclic_sector_family}
With hypotheses as in
Lemma~\ref{thm:exists_common_blocked_cyclic_sector_family}, and given any
integer $p$ that is a common multiple of the per-block period-removal
lengths $(m_k)$, the common reblocking can be performed at the prescribed
length $p$ rather than only at some sufficiently large length.
\end{lemma}

\begin{lemma}[Structural properties of the common reblocked sectors]%
\label{thm:common_blocked_cyclic_sector_derived_properties}
\lean{MPSTensor.CommonBlockedCyclicSectorFamily.derivedProperties}\leanok
\uses{thm:exists_common_blocked_cyclic_sector_family}
Each sector tensor in the common reblocked cyclic-sector family of
Lemma~\ref{thm:exists_common_blocked_cyclic_sector_family} is
trace-preserving ($\sum_i (B^i)^\dagger B^i = \Id$), has primitive
transfer map ($\sigma_\partial(\E_B) = \{1\}$), is tensor-irreducible,
and has positive bond dimension.
\end{lemma}

\begin{lemma}[Common-alphabet flattening]%
\label{thm:common_blocked_cyclic_sector_reindexed_samempv}
\lean{MPSTensor.CommonBlockedCyclicSectorFamily.reindexed_sameMPV}\leanok
\uses{thm:exists_common_blocked_cyclic_sector_family}
With $p = m_k e_k$ for each $k$, the iterated blocked tensor $(B_k^{[m_k]})^{[e_k]}$
agrees with the directly blocked tensor $B_k^{[p]}$ in the same-MPV sense,
after the canonical identification of the $p$-blocked alphabet $d^{[p]}$ with
the $e_k$-fold tensor power of the $m_k$-blocked alphabet $d^{[m_k]}$.
\end{lemma}

\begin{lemma}[Reindexed nonzero-part transport]%
\label{thm:common_blocked_cyclic_sector_reindexed_nonzero_part}
\lean{MPSTensor.CommonBlockedCyclicSectorFamily.reindexed_nonzeroPart}\leanok
\uses{thm:common_blocked_cyclic_sector_reindexed_samempv}
The nonzero-part decomposition of a weighted block sum
$\bigoplus_k \mu_k B_k$ transports through the common-alphabet
identification of
Lemma~\ref{thm:common_blocked_cyclic_sector_reindexed_samempv}: the
common-alphabet sectors have the same nonzero-part decomposition as the
direct $p$-blocked tensors.
\end{lemma}

\begin{definition}[Coordinate-grouping predicate]%
\label{def:common_blocked_cyclic_sector_grouped_block_cast}
\lean{MPSTensor.CommonBlockedCyclicSectorFamily.groupedBlockCastAgrees}\leanok
\uses{thm:exists_common_blocked_cyclic_sector_family}
The canonical identification of the common-alphabet sectors with the iterated
alphabet $(d^{[m_k]})^{e_k}$ agrees with grouping any direct word of length
$p = m_k e_k$ into $e_k$ consecutive words of length $m_k$.
\end{definition}

\begin{lemma}[Digit-level flattening cast]%
\label{thm:common_blocked_cyclic_sector_flatten_word_of_block_cast}
\lean{MPSTensor.CommonBlockedCyclicSectorFamily.flattenWordOfBlock_cast_eq}\leanok
\uses{def:common_blocked_cyclic_sector_grouped_block_cast}
For $p = mn$, the canonical word-flattening cast identifies a $p$-letter
word with the concatenation of $n$ consecutive $m$-letter words: the
two encodings of a word of length $p$ over $d^{[p]}$ and over
$(d^{[m]})^{n}$ are equal under the canonical alphabet identification.
\end{lemma}

\begin{lemma}[Direct vs.\ iterated blocking comparison]%
\label{thm:common_blocked_cyclic_sector_blocked_word_comparison}
\lean{MPSTensor.CommonBlockedCyclicSectorFamily.blocked_word_comparison}\leanok
\uses{thm:common_blocked_cyclic_sector_reindexed_samempv,
      def:common_blocked_cyclic_sector_grouped_block_cast}
Direct blocking at length $p = m_k e_k$ and iterated blocking
$(B_k^{[m_k]})^{[e_k]}$ produce, after the canonical alphabet
identification, the same weighted finite sum: for every $N \ge 0$ and
every blocked word $\sigma$,
\[
  V^{(N)}\!\bigl(B_k^{[p]}\bigr)_\sigma
  \;=\; V^{(N)}\!\bigl((B_k^{[m_k]})^{[e_k]}\bigr)_{\sigma'},
\]
where $\sigma'$ is the image of $\sigma$ under the coordinate-grouping
identification.
\end{lemma}
